% Reference file, not a paper source: WSDM Introduction + Related Work prose (verbatim
% from WSDM_format_revised.tex) with citation placeholders restored from the old AAAI
% draft (pewter_aaai.tex) at the corresponding claims, so real \cite keys can be dropped
% in quickly. Every \citep{PLACEHOLDER: ...} below is copied from pewter_aaai.tex, not
% invented -- see the inline % NOTE comments for the few spots that have no old-version
% counterpart (new WSDM-only content: a self-citation and one still-unwritten result).
\documentclass[sigconf, review, anonymous]{acmart}
\usepackage{amsmath,amssymb,amsthm}
\usepackage{xspace}
\newcommand{\method}{\textsc{Pewter}\xspace}
\newcommand{\Pewter}{\textsc{Pewter}\xspace}

\begin{document}

\section{Introduction}

Many of the graphs that matter on the web put the label on the edge, not on the vertex. A buyer rates a seller as trustworthy or not on a peer-to-peer marketplace, a reviewer marks another reviewer as trusted or distrusted on a product review site, and an editor supports or opposes a candidate in an adminship election. In each case a directed edge carries a sign, most edges are unlabeled or unobserved, and the practical task is to predict the missing signs from the observed ones. Downstream, these predictions drive reputation displays, ranking, and moderation triage, so their failure modes matter as much as their average accuracy.

Other works fold external vertex or edge features, such as review text, timestamps, or account metadata, into the prediction. We do not: an edge's sign is predicted from the graph's topology and the signs already observed on other edges alone.

The classical recipe extracts a fixed set of features around an edge, either external or topological, and classifies the edge from them \citep{PLACEHOLDER: Heider balance theory / Cartwright-Harary structural balance}, \citep{PLACEHOLDER: Leskovec, Huttenlocher, Kleinberg -- signed network sign prediction / status theory}. Such features are hand-built and cannot adapt to the data. To address this, vertex-centric graph neural networks (GNNs) were proposed \citep{PLACEHOLDER: SGCN}, \citep{PLACEHOLDER: SiGAT}, \citep{PLACEHOLDER: SNEA}, \citep{PLACEHOLDER: SDGNN / SGA-GSGNN or other directed signed GNN}: they produce an embedding for each vertex, and classify an edge from the pair of embeddings at its endpoints, optionally with a few edge features. This works when the endpoints determine the label, and it degrades in a way that is predictable once we look at what a vertex embedding can and cannot hold.

A vertex embedding is a single vector that every edge touching that vertex must share. If a vertex sits on edges of several classes, that one vector stands in for all of them at once, and no read-out can recover which class a particular incident edge belongs to. The severity is controlled by the label entropy that survives the coarsening a message-passing GNN applies to vertex identity, and it does not go away with more layers or wider layers, because it is a property of what is being conditioned on, not of the model that conditions. We state this as Proposition~\ref{prop:bottleneck} and, in a capacity form, as Proposition~\ref{prop:capacity}. Note that in a simple graph, each edge is fully defined by the vertices surrounding it, and as such could be claimed to be fully defined. However, for an edge with an unknown class, a pair of vertices with the same WL coloring will induce the same edge class.

We propose an edge-centric approach based on the embedding of an edge and the classification through a random walk on edges. We show that the mutual information between an edge's sign and a neighboring edge's sign collapses by two to nearly five orders of magnitude between line-graph distance 1 and distance 2--3 across all six datasets we study (Section~\ref{sec:empconf}). As such, almost all of the decodable signal about an edge sign sits in the labeled edges one hop away. We thus propose very short walks, which keeps the cost low.

In directed graphs, an edge label need not be equally predictable from its two endpoints. On four of the six datasets the label entropy seen from the source side is significantly lower than from the target side, so the source carries more of the label information; on the two Wikipedia vote graphs the direction significantly reverses. To ensure generality of Pewter, we perform walks starting two edges before the current edge and ending two edges after it, to allow a transformer the choice where to extract the information. We call this directed short-path classifier \Pewter (Proximal Edge-Walk Transformer with Ensemble Read-out).

\section{Related Work}
\paragraph{Edge and sign prediction.} Early work on signed network models predicted edge labels using local graph structure around an edge rather than its endpoints in isolation. Structural balance theory predicts a sign from the parity of a closed triad \citep{PLACEHOLDER: Heider balance theory / Cartwright-Harary structural balance}, and later work replaces balance with status differences and other triad-count features \citep{PLACEHOLDER: Leskovec, Huttenlocher, Kleinberg -- signed network sign prediction / status theory}. These methods rely on a fixed, hand-crafted set of local structural patterns rather than representations learned from data. Pewter keeps the local, edge-centric view of this line of work and replaces the hand-built patterns with a learned sequence model.

\paragraph{Vertex-centric GNNs.} The dominant modern approach replaces hand-built triad features with message-passing GNNs specialized to signed and directed graphs \citep{PLACEHOLDER: SGCN}, \citep{PLACEHOLDER: SiGAT}, \citep{PLACEHOLDER: SNEA}, \citep{PLACEHOLDER: SDGNN / SGA-GSGNN or other directed signed GNN}. These architectures aggregate the signs and directions of incident edges over one or more rounds of message passing, with some variants aggregating in- and out-edges separately to track orientation. Edge labels are then predicted from the resulting pair of endpoint embeddings, optionally concatenated with edge features. Propositions~\ref{prop:bottleneck} and~\ref{prop:capacity} apply to every member of this family, since they constrain the read-out rather than the aggregation.

\paragraph{Walk and sequence representations.} A separate line of work embeds vertices from random walks over the graph \citep{PLACEHOLDER: DeepWalk}, \citep{PLACEHOLDER: node2vec}, including signed-graph extensions of the same idea \citep{PLACEHOLDER: signed/directed walk-based embedding method}. There, the walk generates the training pairs, or context windows, used to fit a per-vertex embedding, and a downstream edge classifier reads a pair of such vertex embeddings, so the vertex bottleneck is reintroduced at the read-out. Sequence models applied directly to walks for vertex or graph classification \citep{PLACEHOLDER: walk-based sequence/RNN or transformer node classifier} process the walk token-by-token, but typically summarize it into one fixed per-vertex or per-graph vector before classification. Pewter differs in that the walk is never summarized into a vertex representation: the prediction is read at the masked edge position itself.

\paragraph{Graph transformers.} Transformer architectures have been adapted to attend over an entire graph or large subgraphs, using structural or positional encodings in place of a fixed adjacency prior \citep{PLACEHOLDER: Graphormer}, \citep{PLACEHOLDER: GPS / SAT or other graph transformer}. Attention allows information to move between distant vertices in a single step rather than only through a fixed number of message-passing hops. These models typically attend over a vertex set and pool the result back down to one representation per vertex before an edge is scored, so they too are covered by Proposition~\ref{prop:bottleneck} whenever the final read-out is a function of two per-vertex vectors. Pewter attends over a short token sequence of alternating vertices and signed edges, which keeps attention cost independent of graph size.

\paragraph{Oversquashing and information bottlenecks.} A related line of work studies how message-passing GNNs lose information as receptive fields grow with depth: each vertex's neighborhood grows exponentially with the number of rounds while the message vector stays a fixed width, so distant information is compressed into a channel too narrow to carry it \citep{PLACEHOLDER: Alon \& Yahav 2021}, \citep{PLACEHOLDER: Topping et al. 2022}. That bottleneck is a property of propagation depth and can in principle be mitigated by rewiring. The bottleneck of Propositions~\ref{prop:bottleneck} and~\ref{prop:capacity} is different in kind: it is a property of the read-out step, and applies to a single fixed-width vertex embedding regardless of how many rounds produced it or how the graph was rewired.

Across these approaches, the mechanism used to propagate sign information differs substantially, but most approaches still predict an edge's label from a pair of fixed per-vertex representations.

\end{document}
